\section {Making your own programs}
\typeout {Section 3 ***}

{\tt LatMRG} offers more flexibility than just providing a set of 
executable programs.
One can also use the modules provided to write one own's programs.  
For that, a Modula-2 compiler and some knowledge about the Modula-2
language are required.
Instead of providing the parameters according to the format of
Appendix A, one sets those parameters (or data) by setting the
appropriate variables or calling the appropriate procedures from
the intermediate or low-level modules.
In this section, we introduce briefly the intermediate and low-level 
modules, which are described in Appendix B and C, respectively.
It is important to recall that the multiplier's and basis components
can be implemented with different representations and that one must make
sure to select the appropriate representation for the target application.
See Section 3.2 for more details on this.

%%%%%%%%%%%%%%%%%%%%%%%%%%%%%%%%%%%%%%%%%%%
\subsection {Using the modules of {\tt LatMRG}}

The modules of {\tt LatMRG} (excluding the executable programs)
have been classified in two sets: lower-level and intermediate-level.
The lower-level modules offer basic facilities for arithmetic operations
and conversions, with different representations, for basis vectors and
multipliers.  They are described in Appendix C.

The intermediate-level module {\tt MRG}
constructs lattice bases for different classes
of generators and perform tests on such generators.
The modules {\tt LATBASIS}, {\tt REDBAS}, {\tt REDBAS2}, and {\tt REDBAS3}
offer tools for manipulating lattice bases,
finding the shortest vector in a lattice, and for reducing
a basis in the sense of Minkowski.
Those modules are described in Appendix B.
The programs of Appendix A use those intermediate and lower-level modules
in their implementation, and so, provide examples of how to use them.

%%%%%%%%%%%%%%%%%%%%%%%%%%%%%%%%%%%%%%%%%%%%%%%%%%%%%
\subsection {Lower-level modules and Changing the representation}

As discussed in Section 1.6, the multiplier's components can be implemented
in the {\tt LONGINT} or {\tt SuperInteger} representation, while the basis
components can be in the {\tt LONGREAL} or {\tt SuperInteger} representation.
To select the appropriate representation (in modules {\tt MULT} and
{\tt BASIS}, respectively), one should use the command
{\tt select} (see Appendix C), set the environment variable LILR,
 and recompile all the modules.

Generally speaking, the proper choice of representation
depends on the size of the mo\-du\-lus $m$. 
For example, if the modulus is less than $2^{31}$, one can type
{\tt select LI LR} (although it could sometimes happen that vector components
get larger than the modulus).
For large moduli, type {\tt select SI SI}.
See the description of the command {\tt select} for more details on its use.
After changing the representation, one must recompile the modules. 
Calling one of the command {\tt lilr, silr, sisi} in a shell will
set the environment variable LILR corresponding to the particular case
of  {\tt select}. (See also the file {\tt Makedoc.tex}.) This is
necessary in order to link with the appropriate libraries.

%%%%%%%%%%%%%%%%%%%%%%%%%%%%%%%%%%%%%%%%%%%%%%%%%%%%%%%%%%%%%%%%
\subsection {Modifying the package, recompiling, and relinking}

After calling {\tt select}, or after modifying the code of a module,
one must recompile and relink. 
\iffalse  Some systems
%%  (e.g., {\em TopSpeed Modula-2\/} on MS-DOS or OS/2)
provide automatic make facilities, which recompile all the modified
modules in the correct order and relink appropriately.
We have written a {\tt Makefile} to facilitate the compiling and
linking of programs which use the lower-level modules of {\tt LatMRG}.
Just type {\tt make} and the name of your program,
and all the modules needed will be recompiled.
See the {\tt Makefile} for more details.
%%  To recompile all the modules except those in executable form,
%%  just type {\tt make all}.
 %%%%%%%%%%%%
There are also two macros, {\tt debug} and {\tt optim} in the makefile,
to add arguments to the compiler and linker commands.
This is useful, for example,
to select the level of object-code optimization:
for full optimization, use the command `{\tt make "optim=-O3"}' and the
name of your program.
\footnote { Est-ce encore valide ? }
\fi  %%%%%%%%%
